\documentclass{amsart}
% For dealing with references we use the comment environment
\usepackage{verbatim}
\newenvironment{reference}{\comment}{\endcomment}
%\newenvironment{reference}{}{}
\newenvironment{slogan}{\comment}{\endcomment}
\newenvironment{history}{\comment}{\endcomment}

% For commutative diagrams we use Xy-pic
\usepackage[all]{xy}

% We use 2cell for 2-commutative diagrams.
\xyoption{2cell}
\UseAllTwocells

% We use multicol for the list of chapters between chapters
\usepackage{multicol}

% This is generally recommended for better output
\usepackage{lmodern}
\usepackage[T1]{fontenc}

% For cross-file-references

% Package for hypertext links:
\usepackage{hyperref}

% For any local file, say "hello.tex" you want to link to please

% Theorem environments.
%
\theoremstyle{plain}
\newtheorem{theorem}[subsection]{Theorem}
\newtheorem{proposition}[subsection]{Proposition}
\newtheorem{lemma}[subsection]{Lemma}

\theoremstyle{definition}
\newtheorem{definition}[subsection]{Definition}
\newtheorem{example}[subsection]{Example}
\newtheorem{exercise}[subsection]{Exercise}
\newtheorem{situation}[subsection]{Situation}

\theoremstyle{remark}
\newtheorem{remark}[subsection]{Remark}
\newtheorem{remarks}[subsection]{Remarks}

\numberwithin{equation}{subsection}

% Macros
%
\def\lim{\mathop{\mathrm{lim}}\nolimits}
\def\colim{\mathop{\mathrm{colim}}\nolimits}
\def\Spec{\mathop{\mathrm{Spec}}}
\def\Hom{\mathop{\mathrm{Hom}}\nolimits}
\def\Ext{\mathop{\mathrm{Ext}}\nolimits}
\def\SheafHom{\mathop{\mathcal{H}\!\mathit{om}}\nolimits}
\def\SheafExt{\mathop{\mathcal{E}\!\mathit{xt}}\nolimits}
\def\Sch{\mathit{Sch}}
\def\Mor{\mathop{\mathrm{Mor}}\nolimits}
\def\Ob{\mathop{\mathrm{Ob}}\nolimits}
\def\Sh{\mathop{\mathit{Sh}}\nolimits}
\def\NL{\mathop{N\!L}\nolimits}
\def\CH{\mathop{\mathrm{CH}}\nolimits}
\def\proetale{{pro\text{-}\acute{e}tale}}
\def\etale{{\acute{e}tale}}
\def\QCoh{\mathit{QCoh}}
\def\Ker{\mathop{\mathrm{Ker}}}
\def\Im{\mathop{\mathrm{Im}}}
\def\Coker{\mathop{\mathrm{Coker}}}
\def\Coim{\mathop{\mathrm{Coim}}}

% Boxtimes
%
\DeclareMathSymbol{\boxtimes}{\mathbin}{AMSa}{"02}

%
% Macros for moduli stacks/spaces
%
\def\QCohstack{\mathcal{QC}\!\mathit{oh}}
\def\Cohstack{\mathcal{C}\!\mathit{oh}}
\def\Spacesstack{\mathcal{S}\!\mathit{paces}}
\def\Quotfunctor{\mathrm{Quot}}
\def\Hilbfunctor{\mathrm{Hilb}}
\def\Curvesstack{\mathcal{C}\!\mathit{urves}}
\def\Polarizedstack{\mathcal{P}\!\mathit{olarized}}
\def\Complexesstack{\mathcal{C}\!\mathit{omplexes}}
% \Pic is the operator that assigns to X its picard group, usage \Pic(X)
% \Picardstack_{X/B} denotes the Picard stack of X over B
% \Picardfunctor_{X/B} denotes the Picard functor of X over B
\def\Pic{\mathop{\mathrm{Pic}}\nolimits}
\def\Picardstack{\mathcal{P}\!\mathit{ic}}
\def\Picardfunctor{\mathrm{Pic}}
\def\Deformationcategory{\mathcal{D}\!\mathit{ef}}

\setlength{\emergencystretch}{3em}
\begin{document}
\title[Stacks Project, Tag 0A9V]{2065 --- Over every ring A, the direct sum of O through O(-n) is a single perfect generator of the unbounded derived category of quasi-coherent modules on projective n-space}
\maketitle
\noindent Copyright (C) 2005--2025 Johan de Jong. Stacks Project authors. Source: \href{https://stacks.math.columbia.edu/tag/0A9V}{Tag 0A9V}.

\noindent Editorial difficulty note (not author text). Difficulty H1: The statement identifies a specific finite vector-bundle generator over an arbitrary base ring and for unbounded quasi-coherent derived objects. The proof constructs every negative twist from one finite Koszul complex and detects an arbitrary nonzero derived object by lifting a local cohomology cycle through ample twists. The unbounded derived-QCoh comparison and finite-generator construction exceed ordinary projective-space cohomology or qualification algebra..

\noindent Editorial provenance note (not author text). History: excerpt from revision \texttt{a0444\allowbreak 6e57e\allowbreak c1fbc\allowbreak 252a8\allowbreak 71afc\allowbreak ec775\allowbreak 2fb28\allowbreak 07b14}, perfect.tex, lines 4154--4193. Reference presentation was adapted; original-author context and core support blocks were retained or added. The mathematical wording is unchanged. Licensed under GNU FDL 1.2 or later; no invariant sections or cover texts. See ../licenses/COPYING. Context blocks reproduce author definitions and supporting results; they are not additional counted problems.

\begin{definition}
\label{derived-definition-generators}
Let $\mathcal{D}$ be a triangulated category. Let $E$ be an object
of $\mathcal{D}$.
\begin{enumerate}
\item We say $E$ is a {\it classical generator} of $\mathcal{D}$
if the smallest strictly full, saturated, triangulated subcategory
of $\mathcal{D}$ containing $E$ is equal to $\mathcal{D}$, in
other words, if $\langle E \rangle = \mathcal{D}$.
\item We say $E$ is a {\it strong generator} of $\mathcal{D}$
if $\langle E \rangle_n = \mathcal{D}$ for some $n \geq 1$.
\item We say $E$ is a {\it weak generator} or a {\it generator}
of $\mathcal{D}$
if for any nonzero object $K$ of $\mathcal{D}$ there exists
an integer $n$ and a nonzero map $E \to K[n]$.
\end{enumerate}
\end{definition}

\begin{lemma}
\label{derived-lemma-right-orthogonal}
Let $\mathcal{D}$ be a triangulated category. Let $E, K$ be objects
of $\mathcal{D}$. The following are equivalent
\begin{enumerate}
\item $\Hom(E, K[i]) = 0$ for all $i \in \mathbf{Z}$,
\item $\Hom(E', K) = 0$ for all $E' \in \langle E \rangle$.
\end{enumerate}
\end{lemma}

\section{The coherator}
\label{section-coherator}

\noindent
Let $X$ be a scheme. The {\it coherator} is a functor
$$
Q_X :
\textit{Mod}(\mathcal{O}_X)
\longrightarrow
\QCoh(\mathcal{O}_X)
$$
which is right adjoint to the inclusion functor
$\QCoh(\mathcal{O}_X) \to \textit{Mod}(\mathcal{O}_X)$.
It exists for any scheme $X$ and moreover the adjunction mapping
$Q_X(\mathcal{F}) \to \mathcal{F}$ is an isomorphism for every
quasi-coherent module $\mathcal{F}$, see
Properties, Proposition \href{https://stacks.math.columbia.edu/tag/077P}{proposition coherator (Tag 077P)}.
Since $Q_X$ is left exact (as a right adjoint) we can consider its
right derived extension
$$
RQ_X :
D(\mathcal{O}_X)
\longrightarrow
D(\QCoh(\mathcal{O}_X)).
$$
Since $Q_X$ is right adjoint to the inclusion functor
$\QCoh(\mathcal{O}_X) \to \textit{Mod}(\mathcal{O}_X)$
we see that $RQ_X$ is right adjoint to the canonical functor
$D(\QCoh(\mathcal{O}_X)) \to D(\mathcal{O}_X)$ by
Derived Categories, Lemma \href{https://stacks.math.columbia.edu/tag/09T5}{lemma derived adjoint functors (Tag 09T5)}.

\medskip\noindent
In this section we will study the functor $RQ_X$. In
Section \href{https://stacks.math.columbia.edu/tag/0CQZ}{section better coherator (Tag 0CQZ)}
we will study the (closely related) right adjoint to the inclusion functor
$D_\QCoh(\mathcal{O}_X) \to D(\mathcal{O}_X)$ (when it exists).



\noindent
In this section we discuss the relationship between quasi-coherent
modules and all modules on a scheme $X$. A reference is
\href{https://stacks.math.columbia.edu/bibliography}{Bibliography: TT (Appendix B)}. By the discussion in
Schemes, Section \href{https://stacks.math.columbia.edu/tag/01LA}{section quasi coherent (Tag 01LA)}
the embedding
$\QCoh(\mathcal{O}_X) \subset \textit{Mod}(\mathcal{O}_X)$
exhibits $\QCoh(\mathcal{O}_X)$ as a weak Serre subcategory of
the category of $\mathcal{O}_X$-modules. Denote
$$
D_\QCoh(\mathcal{O}_X) \subset D(\mathcal{O}_X)
$$
the subcategory of complexes whose cohomology sheaves are quasi-coherent, see
Derived Categories, Section \href{https://stacks.math.columbia.edu/tag/06UP}{section triangulated sub (Tag 06UP)}.
Thus we obtain a canonical functor
\begin{equation}
\label{equation-compare}
D(\QCoh(\mathcal{O}_X))
\longrightarrow
D_\QCoh(\mathcal{O}_X)
\end{equation}
see Derived Categories, Equation (\href{https://stacks.math.columbia.edu/tag/06UR}{equation compare (Tag 06UR)}).

\begin{definition}
\label{definition-supported-on}
Let $X$ be a scheme. Let $E$ be an object of $D(\mathcal{O}_X)$.
Let $T \subset X$ be a closed subset.
We say $E$ is {\it supported on $T$} if the
cohomology sheaves $H^i(E)$ are supported on $T$.
\end{definition}

\begin{definition}
\label{properties-definition-ample}
\begin{reference}
\href{https://stacks.math.columbia.edu/bibliography}{Bibliography: EGA (II Definition 4.5.3)}
\end{reference}
Let $X$ be a scheme.
Let $\mathcal{L}$ be an invertible $\mathcal{O}_X$-module.
We say $\mathcal{L}$ is {\it ample} if
\begin{enumerate}
\item $X$ is quasi-compact, and
\item for every $x \in X$ there exists an $n \geq 1$
and $s \in \Gamma(X, \mathcal{L}^{\otimes n})$ such
that $x \in X_s$ and $X_s$ is affine.
\end{enumerate}
\end{definition}

\begin{lemma}
\label{properties-lemma-open-in-proj-ample}
Let $X$ be a scheme.
Let $S$ be a graded ring. Assume $X$ is quasi-compact,
and assume there exists an open immersion
$$
j : X \longrightarrow Y = \text{Proj}(S).
$$
Then $j^*\mathcal{O}_Y(d)$ is an invertible ample sheaf
for some $d > 0$.
\end{lemma}

\begin{definition}
\label{constructions-definition-projective-space}
The scheme
$\mathbf{P}^n_{\mathbf{Z}} = \text{Proj}(\mathbf{Z}[T_0, \ldots, T_n])$
is called {\it projective $n$-space over $\mathbf{Z}$}.
Its base change $\mathbf{P}^n_S$ to a scheme $S$ is called
{\it projective $n$-space over $S$}. If $R$ is a ring the base change
to $\Spec(R)$ is denoted $\mathbf{P}^n_R$ and called
{\it projective $n$-space over $R$}.
\end{definition}

\begin{definition}
\label{more-algebra-definition-koszul-complex}
Let $R$ be a ring and let $f_1, \ldots, f_r \in R$. The
{\it Koszul complex on $f_1, \ldots, f_r$} is the Koszul complex
associated to the map $(f_1, \ldots, f_r) : R^{\oplus r} \to R$.
Notation $K_\bullet(f_\bullet)$, $K_\bullet(f_1, \ldots, f_r)$,
$K_\bullet(R, f_1, \ldots, f_r)$, or $K_\bullet(R, f_\bullet)$.
\end{definition}

\begin{definition}
\label{modules-definition-gamma-star}
Let $(X, \mathcal{O}_X)$ be a ringed space.
Given an invertible sheaf $\mathcal{L}$ on $X$ we define
the {\it associated graded ring} to be
$$
\Gamma_*(X, \mathcal{L})
=
\bigoplus\nolimits_{n \geq 0} \Gamma(X, \mathcal{L}^{\otimes n})
$$
Given a sheaf of $\mathcal{O}_X$-modules $\mathcal{F}$ we set
$$
\Gamma_*(X, \mathcal{L}, \mathcal{F})
=
\bigoplus\nolimits_{n \in \mathbf{Z}} \Gamma(X,
\mathcal{F} \otimes_{\mathcal{O}_X} \mathcal{L}^{\otimes n})
$$
which we think of as a graded $\Gamma_*(X, \mathcal{L})$-module.
\end{definition}

\begin{lemma}
\label{lemma-nonzero-some-cohomology}
Let $X$ be a scheme and $\mathcal{L}$ an ample invertible
$\mathcal{O}_X$-module. If $K$ is a nonzero object of
$D_\QCoh(\mathcal{O}_X)$, then for some $n \geq 0$ and $p \in \mathbf{Z}$
the cohomology group
$H^p(X, K \otimes_{\mathcal{O}_X}^\mathbf{L} \mathcal{L}^{\otimes n})$
is nonzero.
\end{lemma}

\begin{proof}
Recall that as $X$ has an ample invertible sheaf, it is quasi-compact
and separated (Properties, Definition \ref{properties-definition-ample} and
Lemma \href{https://stacks.math.columbia.edu/tag/01PY}{lemma affine s opens cover quasi separated (Tag 01PY)}).
Thus we may apply
Proposition \ref{proposition-quasi-compact-affine-diagonal}
and represent $K$ by a complex $\mathcal{F}^\bullet$ of
quasi-coherent modules. Pick any $p$ such that
$\mathcal{H}^p = \Ker(\mathcal{F}^p \to \mathcal{F}^{p + 1})/
\Im(\mathcal{F}^{p - 1} \to \mathcal{F}^p)$ is nonzero.
Choose a point $x \in X$ such that the stalk $\mathcal{H}^p_x$ is
nonzero. Choose an $n \geq 0$ and $s \in \Gamma(X, \mathcal{L}^{\otimes n})$
such that $X_s$ is an affine open neighbourhood of $x$.
Choose $\tau \in \mathcal{H}^p(X_s)$ which maps to a nonzero
element of the stalk $\mathcal{H}^p_x$; this is possible
as $\mathcal{H}^p$ is quasi-coherent and $X_s$ is affine.
Since taking sections over $X_s$ is an exact functor on
quasi-coherent modules, we can find a section $\tau' \in \mathcal{F}^p(X_s)$
mapping to zero in $\mathcal{F}^{p + 1}(X_s)$ and mapping to
$\tau$ in $\mathcal{H}^p(X_s)$. By
Properties, Lemma \ref{properties-lemma-invert-s-sections}
there exists an $m$ such that $\tau' \otimes s^{\otimes m}$
is the image of a section
$\tau'' \in \Gamma(X, \mathcal{F}^p \otimes \mathcal{L}^{\otimes mn})$.
Applying the same lemma once more, we find $l \geq 0$ such that
$\tau'' \otimes s^{\otimes l}$ maps to zero in
$\mathcal{F}^{p + 1} \otimes \mathcal{L}^{\otimes (m + l)n}$.
Then $\tau''$ gives a nonzero class in
$H^p(X, K \otimes^\mathbf{L}_{\mathcal{O}_X} \mathcal{L}^{(m + l)n})$
as desired.
\end{proof}


\begin{lemma}
\label{lemma-construct-the-next-one}
Let $A$ be a ring. Let $X = \mathbf{P}^n_A$. For every $a \in \mathbf{Z}$
there exists an exact complex
$$
0 \to \mathcal{O}_X(a) \to \ldots
\to \mathcal{O}_X(a + i)^{\oplus {n + 1 \choose i}} \to
\ldots \to \mathcal{O}_X(a + n + 1) \to 0
$$
of vector bundles on $X$.
\end{lemma}

\begin{proof}
Recall that $\mathbf{P}^n_A$ is $\text{Proj}(A[X_0, \ldots, X_n])$, see
Constructions, Definition \ref{constructions-definition-projective-space}.
Consider the Koszul complex
$$
K_\bullet = K_\bullet(A[X_0, \ldots, X_n], X_0, \ldots, X_n)
$$
over $S = A[X_0, \ldots, X_n]$ on $X_0, \ldots, X_n$.
Since $X_0, \ldots, X_n$ is clearly a regular sequence in the
polynomial ring $S$, we see that
(More on Algebra, Lemma \ref{more-algebra-lemma-regular-koszul-regular})
that the Koszul complex $K_\bullet$ is exact, except in degree $0$
where the cohomology is $S/(X_0, \ldots, X_n)$.
Note that $K_\bullet$ becomes a complex of graded modules if we
put the generators of $K_i$ in degree $+i$. In other words an
exact complex
$$
0 \to S(-n - 1) \to \ldots \to S(-n - 1 + i)^{\oplus {n \choose i}} \to \ldots
\to S \to S/(X_0, \ldots, X_n) \to 0
$$
Applying the exact functor $\tilde{\ }$ functor of Constructions, 
Lemma \ref{constructions-lemma-proj-sheaves} and using that
the last term is in the kernel of this functor,
we obtain the exact complex
$$
0 \to \mathcal{O}_X(-n - 1) \to \ldots
\to \mathcal{O}_X(-n - 1 + i)^{\oplus {n + 1 \choose i}} \to
\ldots \to \mathcal{O}_X \to 0
$$
Twisting by the invertible sheaves $\mathcal{O}_X(n + a + 1)$
we get the exact complexes of the lemma.
\end{proof}


\begin{lemma}
\label{lemma-argument-proves}
Let $X$ be a quasi-compact and quasi-separated scheme. Suppose that
for every affine open $U \subset X$ the right derived functor
$$
\Phi : D(\QCoh(\mathcal{O}_U)) \to D(\QCoh(\mathcal{O}_X))
$$
of the left exact functor
$j_* : \QCoh(\mathcal{O}_U) \to \QCoh(\mathcal{O}_X)$
fits into a commutative diagram
$$
\xymatrix{
D(\QCoh(\mathcal{O}_U)) \ar[d]_\Phi \ar[r]_{i_U} &
D_\QCoh(\mathcal{O}_U) \ar[d]^{Rj_*} \\
D(\QCoh(\mathcal{O}_X)) \ar[r]^{i_X} &
D_\QCoh(\mathcal{O}_X)
}
$$
Then the functor (\ref{equation-compare})
$$
D(\QCoh(\mathcal{O}_X))
\longrightarrow
D_\QCoh(\mathcal{O}_X)
$$
is an equivalence with quasi-inverse given by $RQ_X$.
\end{lemma}

\begin{proof}
Let $E$ be an object of $D_\QCoh(\mathcal{O}_X)$ and
let $A$ be an object of $D(\QCoh(\mathcal{O}_X))$.
We have to show that the adjunction maps
$$
A \to RQ_X(i_X(A))
\quad\text{and}\quad
i_X(RQ_X(E)) \to E
$$
are isomorphisms. Consider the hypothesis
$H_n$: the adjunction maps above are isomorphisms
whenever $E$ and $i_X(A)$ are supported
(Definition \ref{definition-supported-on})
on a closed subset of $X$ which
is contained in the union of $n$ affine opens of $X$.
We will prove $H_n$ by induction on $n$.

\medskip\noindent
Base case: $n = 0$. In this case $E = 0$, hence the map
$i_X(RQ_X(E)) \to E$ is an isomorphism. Similarly $i_X(A) = 0$.
Thus the cohomology sheaves of $i_X(A)$ are zero. Since the inclusion
functor $\QCoh(\mathcal{O}_X) \to \textit{Mod}(\mathcal{O}_X)$
is fully faithful and exact, we conclude that the cohomology
objects of $A$ are zero, i.e., $A = 0$ and
$A \to RQ_X(i_X(A))$ is an isomorphism as well.

\medskip\noindent
Induction step. Suppose that $E$ and $i_X(A)$ are supported on a
closed subset $T$ of $X$ contained in $U_1 \cup \ldots \cup U_n$
with $U_i \subset X$ affine open. Set $U = U_n$.
Consider the distinguished triangles
$$
A \to \Phi(A|_U) \to A' \to A[1]
\quad\text{and}\quad
E \to Rj_*(E|_U) \to E' \to E[1]
$$
where $\Phi$ is as in the statement of the lemma.
Note that $E \to Rj_*(E|_U)$ is a quasi-isomorphism over $U = U_n$.
Since $i_X \circ \Phi = Rj_* \circ i_U$ by assumption
and since $i_X(A)|_U = i_U(A|_U)$
we see that $i_X(A) \to i_X(\Phi(A|_U))$ is a quasi-isomorphism over $U$.
Hence $i_X(A')$ and $E'$ are supported on the closed
subset $T \setminus U$ of $X$ which is contained in
$U_1 \cup \ldots \cup U_{n - 1}$.
By induction hypothesis the statement is true for $A'$ and $E'$. By
Derived Categories, Lemma \href{https://stacks.math.columbia.edu/tag/014A}{lemma third isomorphism triangle (Tag 014A)}
it suffices to prove the maps
$$
\Phi(A|_U) \to RQ_X(i_X(\Phi(A|_U)))
\quad\text{and}\quad
i_X(RQ_X(Rj_*E|_U)) \to Rj_*(E|_U)
$$
are isomorphisms. By assumption and by
Lemma \ref{lemma-flat-pushforward-coherator}
(the inclusion morphism $j : U \to X$ is flat, quasi-compact, and
quasi-separated) we have
$$
RQ_X(i_X(\Phi(A|_U))) = RQ_X(Rj_*(i_U(A|_U))) = \Phi(RQ_U(i_U(A|_U)))
$$
and
$$
i_X(RQ_X(Rj_*(E|_U))) = i_X(\Phi(RQ_U(E|_U))) = Rj_*(i_U(RQ_U(E|_U)))
$$
Finally, the maps
$$
A|_U \to RQ_U(i_U(A|_U))
\quad\text{and}\quad
i_U(RQ_U(E|_U)) \to E|_U
$$
are isomorphisms by Lemma \ref{lemma-affine-coherator}. The result follows.
\end{proof}


\begin{lemma}
\label{lemma-affine-pushforward}
Let $f : X \to Y$ be an affine morphism of schemes.
Then $f_*$ defines a derived functor
$f_* : D(\QCoh(\mathcal{O}_X)) \to D(\QCoh(\mathcal{O}_Y))$.
This functor has the property that
$$
\xymatrix{
D(\QCoh(\mathcal{O}_X)) \ar[d]_{f_*} \ar[r] &
D_\QCoh(\mathcal{O}_X) \ar[d]^{Rf_*} \\
D(\QCoh(\mathcal{O}_Y)) \ar[r] &
D_\QCoh(\mathcal{O}_Y)
}
$$
commutes.
\end{lemma}

\begin{proof}
The functor
$f_* : \QCoh(\mathcal{O}_X) \to \QCoh(\mathcal{O}_Y)$
is exact, see
Cohomology of Schemes, Lemma \href{https://stacks.math.columbia.edu/tag/01XC}{lemma relative affine vanishing (Tag 01XC)}.
Hence $f_*$ defines a derived functor
$f_* : D(\QCoh(\mathcal{O}_X)) \to D(\QCoh(\mathcal{O}_Y))$
by simply applying $f_*$ to any representative complex, see
Derived Categories, Lemma \href{https://stacks.math.columbia.edu/tag/015F}{lemma right derived exact functor (Tag 015F)}.
The diagram commutes by Lemma \ref{lemma-pushforward-affine-morphism}.
\end{proof}


\begin{lemma}
\label{lemma-flat-pushforward-coherator}
Let $f : X \to Y$ be a morphism of schemes. Assume $f$ is
quasi-compact, quasi-separated, and flat. Then, denoting
$$
\Phi : D(\QCoh(\mathcal{O}_X)) \to D(\QCoh(\mathcal{O}_Y))
$$
the right derived functor of
$f_* : \QCoh(\mathcal{O}_X) \to \QCoh(\mathcal{O}_Y)$
we have $RQ_Y \circ Rf_* = \Phi \circ RQ_X$.
\end{lemma}

\begin{proof}
We will prove this by showing that $RQ_Y \circ Rf_*$ and $\Phi \circ RQ_X$
are right adjoint to the same functor
$D(\QCoh(\mathcal{O}_Y)) \to D(\mathcal{O}_X)$.

\medskip\noindent
Since $f$ is quasi-compact and quasi-separated, we see that
$f_*$ preserves quasi-coherence, see
Schemes, Lemma \href{https://stacks.math.columbia.edu/tag/01LC}{lemma push forward quasi coherent (Tag 01LC)}.
Recall that $\QCoh(\mathcal{O}_X)$ is a Grothendieck abelian category
(Properties, Proposition \href{https://stacks.math.columbia.edu/tag/077P}{proposition coherator (Tag 077P)}).
Hence any $K$ in $D(\QCoh(\mathcal{O}_X))$
can be represented by a K-injective complex $\mathcal{I}^\bullet$
of $\QCoh(\mathcal{O}_X)$, see
Injectives, Theorem
\href{https://stacks.math.columbia.edu/tag/079P}{theorem K injective embedding grothendieck (Tag 079P)}.
Then we can define $\Phi(K) = f_*\mathcal{I}^\bullet$.

\medskip\noindent
Since $f$ is flat, the functor $f^*$ is exact. Hence $f^*$ defines
$f^* : D(\mathcal{O}_Y) \to D(\mathcal{O}_X)$ and also
$f^* : D(\QCoh(\mathcal{O}_Y)) \to D(\QCoh(\mathcal{O}_X))$.
The functor $f^* = Lf^* : D(\mathcal{O}_Y) \to D(\mathcal{O}_X)$
is left adjoint to
$Rf_* : D(\mathcal{O}_X) \to D(\mathcal{O}_Y)$,
see Cohomology, Lemma \href{https://stacks.math.columbia.edu/tag/079W}{lemma adjoint (Tag 079W)}.
Similarly, the functor
$f^* : D(\QCoh(\mathcal{O}_Y)) \to D(\QCoh(\mathcal{O}_X))$
is left adjoint to
$\Phi : D(\QCoh(\mathcal{O}_X)) \to D(\QCoh(\mathcal{O}_Y))$
by Derived Categories, Lemma \href{https://stacks.math.columbia.edu/tag/09T5}{lemma derived adjoint functors (Tag 09T5)}.

\medskip\noindent
Let $A$ be an object of $D(\QCoh(\mathcal{O}_Y))$ and
$E$ an object of $D(\mathcal{O}_X)$. Then
\begin{align*}
\Hom_{D(\QCoh(\mathcal{O}_Y))}(A, RQ_Y(Rf_*E))
& =
\Hom_{D(\mathcal{O}_Y)}(A, Rf_*E) \\
& =
\Hom_{D(\mathcal{O}_X)}(f^*A, E) \\
& =
\Hom_{D(\QCoh(\mathcal{O}_X))}(f^*A, RQ_X(E)) \\
& =
\Hom_{D(\QCoh(\mathcal{O}_Y))}(A, \Phi(RQ_X(E)))
\end{align*}
This implies what we want.
\end{proof}


\begin{lemma}
\label{lemma-affine-coherator}
Let $X = \Spec(A)$ be an affine scheme. Then
\begin{enumerate}
\item $Q_X : \textit{Mod}(\mathcal{O}_X) \to \QCoh(\mathcal{O}_X)$
is the functor
which sends $\mathcal{F}$ to the quasi-coherent $\mathcal{O}_X$-module
associated to the $A$-module $\Gamma(X, \mathcal{F})$,
\item $RQ_X : D(\mathcal{O}_X) \to D(\QCoh(\mathcal{O}_X))$
is the functor which sends $E$ to the complex of quasi-coherent
$\mathcal{O}_X$-modules associated to the object $R\Gamma(X, E)$ of $D(A)$,
\item restricted to $D_\QCoh(\mathcal{O}_X)$ the functor
$RQ_X$ defines a quasi-inverse to (\ref{equation-compare}).
\end{enumerate}
\end{lemma}

\begin{proof}
The functor $Q_X$ is the functor
$$
\mathcal{F} \mapsto \widetilde{\Gamma(X, \mathcal{F})}
$$
by Schemes, Lemma \href{https://stacks.math.columbia.edu/tag/01I7}{lemma compare constructions (Tag 01I7)}.
This immediately implies (1) and (2). The third assertion
follows from (the proof of)
Lemma \ref{lemma-affine-compare-bounded}.
\end{proof}


\begin{lemma}
\label{lemma-affine-compare-bounded}
Let $X = \Spec(A)$ be an affine scheme. All the functors in the diagram
$$
\xymatrix{
D(\QCoh(\mathcal{O}_X)) \ar[rr]_{(\ref{equation-compare})}
& &
D_\QCoh(\mathcal{O}_X) \ar[ld]^{R\Gamma(X, -)} \\
& D(A) \ar[lu]^{\widetilde{\ \ }}
}
$$
are equivalences of triangulated categories. Moreover, for $E$ in
$D_\QCoh(\mathcal{O}_X)$ we have $H^0(X, E) = H^0(X, H^0(E))$.
\end{lemma}

\begin{proof}
The functor $R\Gamma(X, -)$ gives a functor
$D(\mathcal{O}_X) \to D(A)$ and hence by restriction a functor
\begin{equation}
\label{equation-back}
R\Gamma(X, -) : D_\QCoh(\mathcal{O}_X) \longrightarrow D(A).
\end{equation}
We will show this functor is quasi-inverse to (\ref{equation-compare})
via the equivalence between quasi-coherent modules on $X$ and
the category of $A$-modules.

\medskip\noindent
Elucidation. Denote $(Y, \mathcal{O}_Y)$ the one point space with sheaf
of rings given by $A$. Denote
$\pi : (X, \mathcal{O}_X) \to (Y, \mathcal{O}_Y)$
the obvious morphism of ringed spaces.
Then $R\Gamma(X, -)$ can be identified with $R\pi_*$ and the functor
(\ref{equation-compare}) via the equivalence
$\textit{Mod}(\mathcal{O}_Y) = \text{Mod}_A = \QCoh(\mathcal{O}_X)$
can be identified with $L\pi^* = \pi^* = \widetilde{\ }$ (see
Modules, Lemma \href{https://stacks.math.columbia.edu/tag/01BH}{lemma construct quasi coherent sheaves (Tag 01BH)} and
Schemes, Lemmas \href{https://stacks.math.columbia.edu/tag/01I7}{lemma compare constructions (Tag 01I7)} and
\href{https://stacks.math.columbia.edu/tag/01IB}{lemma equivalence quasi coherent (Tag 01IB)}). Thus the functors
$$
\xymatrix{
D(A) \ar@<1ex>[r] & D(\mathcal{O}_X) \ar@<1ex>[l]
}
$$
are adjoint (by Cohomology, Lemma \href{https://stacks.math.columbia.edu/tag/079W}{lemma adjoint (Tag 079W)}). In
particular we obtain canonical adjunction mappings
$$
a : \widetilde{R\Gamma(X, E)} \longrightarrow E
$$
for $E$ in $D(\mathcal{O}_X)$ and
$$
b : M^\bullet \longrightarrow R\Gamma(X, \widetilde{M^\bullet})
$$
for $M^\bullet$ a complex of $A$-modules.

\medskip\noindent
Let $E$ be an object of $D_\QCoh(\mathcal{O}_X)$. We may apply
Lemma \href{https://stacks.math.columbia.edu/tag/08D4}{lemma application nice K injective (Tag 08D4)}
to the functor $F(-) = \Gamma(X, -)$
with $N = 1$ by Cohomology of Schemes, Lemma
\href{https://stacks.math.columbia.edu/tag/01XB}{lemma quasi coherent affine cohomology zero (Tag 01XB)}.
Hence
$$
H^0(R\Gamma(X, E)) = H^0(R\Gamma(X, \tau_{\geq 0}E)) = \Gamma(X, H^0(E))
$$
(the last equality by definition of the canonical truncation).
Using this we will show that the adjunction mappings $a$ and $b$
induce isomorphisms $H^0(a)$ and $H^0(b)$. Thus $a$ and $b$
are quasi-isomorphisms (as the statement is invariant under shifts)
and the lemma is proved.

\medskip\noindent
In both cases we use that $\widetilde{\ }$ is an exact functor
(Schemes, Lemma \href{https://stacks.math.columbia.edu/tag/01HV}{lemma spec sheaves (Tag 01HV)}). Namely, this
implies that
$$
H^0\left(\widetilde{R\Gamma(X, E)}\right) =
\widetilde{H^0(R\Gamma(X, E))} =
\widetilde{\Gamma(X, H^0(E))}
$$
which is equal to $H^0(E)$ because $H^0(E)$ is quasi-coherent. Thus
$H^0(a)$ is an isomorphism. For the other direction we have
$$
H^0(R\Gamma(X, \widetilde{M^\bullet})) =
\Gamma(X, H^0(\widetilde{M^\bullet})) =
\Gamma(X, \widetilde{H^0(M^\bullet)}) =
H^0(M^\bullet)
$$
which proves that $H^0(b)$ is an isomorphism.
\end{proof}


\begin{lemma}
\label{lemma-quasi-coherence-pushforward}
If $f : X \to Y$ is a morphism of affine schemes given by the ring map
$A \to B$, then the diagram
$$
\xymatrix{
D(B) \ar[d] \ar[r] & D_\QCoh(\mathcal{O}_X) \ar[d]^{Rf_*} \\
D(A) \ar[r] & D_\QCoh(\mathcal{O}_Y)
}
$$
commutes.
\end{lemma}

\begin{proof}
Follows from Lemma \ref{lemma-affine-compare-bounded}
using that $R\Gamma(Y, Rf_*K) = R\Gamma(X, K)$ by
Cohomology, Lemma \href{https://stacks.math.columbia.edu/tag/0D5W}{lemma Leray unbounded (Tag 0D5W)}.
\end{proof}


\begin{lemma}
\label{lemma-quasi-coherence-direct-image}
Let $f : X \to S$ be a morphism of schemes.
Assume that $f$ is quasi-separated and quasi-compact.
\begin{enumerate}
\item The functor $Rf_*$ sends $D_\QCoh(\mathcal{O}_X)$
into $D_\QCoh(\mathcal{O}_S)$.
\item If $S$ is quasi-compact, there exists an integer $N = N(X, S, f)$
such that for an object $E$ of $D_\QCoh(\mathcal{O}_X)$
with $H^m(E) = 0$ for $m > 0$ we have
$H^m(Rf_*E) = 0$ for $m \geq N$.
\item In fact, if $S$ is quasi-compact we can find $N = N(X, S, f)$
such that for every morphism of schemes $S' \to S$
the same conclusion holds for the functor $R(f')_*$
where $f' : X' \to S'$ is the base change of $f$.
\end{enumerate}
\end{lemma}

\begin{proof}
Let $E$ be an object of $D_\QCoh(\mathcal{O}_X)$. To prove (1) we have to
show that $Rf_*E$ has quasi-coherent cohomology sheaves. The question is local
on $S$, hence we may assume $S$ is quasi-compact. Pick $N = N(X, S, f)$ as in
Cohomology of Schemes, Lemma
\href{https://stacks.math.columbia.edu/tag/01XJ}{lemma quasi coherence higher direct images (Tag 01XJ)}.
Thus $R^pf_*\mathcal{F} = 0$ for all quasi-coherent $\mathcal{O}_X$-modules
$\mathcal{F}$ and all $p \geq N$ and the same remains true after base change.

\medskip\noindent
First, assume $E$ is bounded below. We will show (1) and (2) and (3) hold
for such $E$ with our choice of $N$. In this case we can for example use the
spectral sequence
$$
R^pf_*H^q(E) \Rightarrow R^{p + q}f_*E
$$
(Derived Categories, Lemma \href{https://stacks.math.columbia.edu/tag/015J}{lemma two ss complex functor (Tag 015J)}),
the quasi-coherence of $R^pf_*H^q(E)$, and the vanishing of $R^pf_*H^q(E)$
for $p \geq N$ to see that (1), (2), and (3) hold in this case.

\medskip\noindent
Next we prove (2) and (3). Say $H^m(E) = 0$ for $m > 0$.
Let $U \subset S$ be affine open. By Cohomology of Schemes, Lemma
\href{https://stacks.math.columbia.edu/tag/01XK}{lemma quasi coherence higher direct images application (Tag 01XK)}
and our choice of $N$
we have $H^p(f^{-1}(U), \mathcal{F}) = 0$ for $p \geq N$
and any quasi-coherent $\mathcal{O}_X$-module $\mathcal{F}$.
Hence we may apply Lemma \href{https://stacks.math.columbia.edu/tag/08D4}{lemma application nice K injective (Tag 08D4)}
to the functor $\Gamma(f^{-1}(U), -)$ to see that
$$
R\Gamma(U, Rf_*E) = R\Gamma(f^{-1}(U), E)
$$
has vanishing cohomology in degrees $\geq N$. Since this holds for
all $U \subset S$ affine open we conclude that $H^m(Rf_*E) = 0$
for $m \geq N$.

\medskip\noindent
Next, we prove (1) in the general case. Recall that there is a
distinguished triangle
$$
\tau_{\leq -n - 1}E \to E \to \tau_{\geq -n}E \to
(\tau_{\leq -n - 1}E)[1]
$$
in $D(\mathcal{O}_X)$, see Derived Categories, Remark
\href{https://stacks.math.columbia.edu/tag/08J5}{remark truncation distinguished triangle (Tag 08J5)}.
By (2) we see that $Rf_*\tau_{\leq -n - 1}E$
has vanishing cohomology sheaves in degrees $\geq -n + N$.
Thus, given an integer $q$ we see that $R^qf_*E$ is equal
to $R^qf_*\tau_{\geq -n}E$ for some $n$ and the result
above applies.
\end{proof}


\begin{lemma}
\label{lemma-acyclicity-lemma}
Let $f : X \to S$ be a quasi-separated and quasi-compact morphism
of schemes. Let $\mathcal{F}^\bullet$ be a complex of quasi-coherent
$\mathcal{O}_X$-modules each of which is right acyclic for $f_*$.
Then $f_*\mathcal{F}^\bullet$ represents $Rf_*\mathcal{F}^\bullet$
in $D(\mathcal{O}_S)$.
\end{lemma}

\begin{proof}
There is always a canonical map
$f_*\mathcal{F}^\bullet \to Rf_*\mathcal{F}^\bullet$.
Our task is to show that this is an isomorphism on cohomology sheaves.
As the statement is invariant under shifts it suffices to show that
$H^0(f_*(\mathcal{F}^\bullet)) \to R^0f_*\mathcal{F}^\bullet$
is an isomorphism. The statement is local on $S$ hence we
may assume $S$ affine. By
Lemma \ref{lemma-quasi-coherence-direct-image}
we have $R^0f_*\mathcal{F}^\bullet = R^0f_*\tau_{\geq -n}\mathcal{F}^\bullet$
for all sufficiently large $n$. Thus we may assume $\mathcal{F}^\bullet$
bounded below. As each $\mathcal{F}^n$ is right $f_*$-acyclic by
assumption we see that $f_*\mathcal{F}^\bullet \to Rf_*\mathcal{F}^\bullet$
is a quasi-isomorphism by Leray's acyclicity lemma (Derived Categories, Lemma
\href{https://stacks.math.columbia.edu/tag/015E}{lemma leray acyclicity (Tag 015E)}).
\end{proof}


\begin{lemma}
\label{lemma-pushforward-affine-morphism}
Let $f : X \to S$ be an affine morphism of schemes. Let $\mathcal{F}^\bullet$
be a complex of quasi-coherent $\mathcal{O}_X$-modules. Then
$f_*\mathcal{F}^\bullet = Rf_*\mathcal{F}^\bullet$.
\end{lemma}

\begin{proof}
Combine Lemma \ref{lemma-acyclicity-lemma} with
Cohomology of Schemes, Lemma \href{https://stacks.math.columbia.edu/tag/01XC}{lemma relative affine vanishing (Tag 01XC)}.
An alternative proof is to work affine locally on $S$
and use Lemma \ref{lemma-quasi-coherence-pushforward}.
\end{proof}


\begin{lemma}
\label{properties-lemma-invert-s-sections}
\begin{reference}
\href{https://stacks.math.columbia.edu/bibliography}{Bibliography: EGA1-second (Section 6.8)}
\end{reference}
Let $X$ be a scheme. Let $\mathcal{L}$ be an invertible sheaf on $X$.
Let $s \in \Gamma(X, \mathcal{L})$. Let $\mathcal{F}$ be a quasi-coherent
$\mathcal{O}_X$-module.
\begin{enumerate}
\item If $X$ is quasi-compact, then (\ref{properties-equation-module-invert-s})
is injective, and
\item if $X$ is quasi-compact and quasi-separated, then
(\ref{properties-equation-module-invert-s}) is an isomorphism.
\end{enumerate}
In particular, the canonical map
$$
\Gamma_*(X, \mathcal{L})_{(s)}
\longrightarrow
\Gamma(X_s, \mathcal{O}_X),\quad
a/s^n \longmapsto a \otimes s^{-n}
$$
is an isomorphism if $X$ is quasi-compact and quasi-separated.
\end{lemma}

\begin{proof}
Assume $X$ is quasi-compact. Choose a finite affine open covering
$X = U_1 \cup \ldots \cup U_m$ with $U_j$ affine and
$\mathcal{L}|_{U_j} \cong \mathcal{O}_{U_j}$. Via this isomorphism,
the image $s|_{U_j}$ corresponds to some
$f_j \in \Gamma(U_j, \mathcal{O}_{U_j})$. Then
$X_s \cap U_j = D(f_j)$.

\medskip\noindent
Proof of (1). Let $t/s^n$ be an element in the kernel of
(\ref{properties-equation-module-invert-s}). Then $t|_{X_s} = 0$.
Hence $(t|_{U_j})|_{D(f_j)} = 0$. By
Lemma \ref{properties-lemma-invert-f-sections} we conclude that
$f_j^{e_j} t|_{U_j} = 0$ for some
$e_j \geq 0$. Let $e = \max(e_j)$. Then we see that $t \otimes s^e$
restricts to zero on $U_j$ for all $j$, hence is zero. Since $t/s^n$
is equal to $t \otimes s^e/s^{n + e}$ in
$\Gamma_*(X, \mathcal{L}, \mathcal{F})_{(s)}$ we conclude that $t/s^n = 0$
as desired.

\medskip\noindent
Proof of (2). Assume $X$ is quasi-compact and quasi-separated.
Then $U_j \cap U_{j'}$ is quasi-compact for all pairs $j, j'$, see
Schemes, Lemma \href{https://stacks.math.columbia.edu/tag/01KO}{lemma characterize quasi separated (Tag 01KO)}.
By part (1) we know (\ref{properties-equation-module-invert-s}) is injective.
Let $t' \in \Gamma(X_s, \mathcal{F}|_{X_s})$. For every $j$, there exist an
integer $e_j \geq 0$ and $t'_j \in \Gamma(U_j, \mathcal{F}|_{U_j})$ such that
$t'|_{D(f_j)}$ corresponds to $t'_j/f_j^{e_j}$
via the isomorphism of Lemma \ref{properties-lemma-invert-f-sections}.
Set $e = \max(e_j)$ and
$$
t_j = f_j^{e - e_j} t'_j \otimes q_j^e \in
\Gamma(U_j,
(\mathcal{F} \otimes_{\mathcal{O}_X} \mathcal{L}^{\otimes e})|_{U_j})
$$
where $q_j \in \Gamma(U_j, \mathcal{L}|_{U_j})$ is the trivializing
section coming from the isomorphism
$\mathcal{L}|_{U_j} \cong \mathcal{O}_{U_j}$. In particular we have
$s|_{U_j} = f_j q_j$. Using this a calculation shows that
$t_j|_{U_j \cap U_{j'}}$ and $t_{j'}|_{U_j \cap U_{j'}}$
map to the same section of $\mathcal{F}$ over $U_j \cap U_{j'} \cap X_s$.
By quasi-compactness of $U_j \cap U_{j'}$ and part (1) there exists an
integer $e' \geq 0$ such that
$$
t_j|_{U_j \cap U_{j'}} \otimes s^{e'}|_{U_j \cap U_{j'}} =
t_{j'}|_{U_j \cap U_{j'}} \otimes s^{e'}|_{U_j \cap U_{j'}}
$$
as sections of $\mathcal{F} \otimes \mathcal{L}^{\otimes e + e'}$ over
$U_j \cap U_{j'}$. We may choose the same $e'$ to work for all pairs
$j, j'$. Then the sheaf conditions implies there is a section
$t \in \Gamma(X, \mathcal{F} \otimes \mathcal{L}^{\otimes e + e'})$
whose restriction to $U_j$ is $t_j \otimes s^{e'}|_{U_j}$.
A simple computation shows that $t/s^{e + e'}$ maps to $t'$
as desired.
\end{proof}


\begin{lemma}
\label{properties-lemma-invert-f-sections}
\begin{slogan}
Sections of quasi-coherent sheaves have only meromorphic singularities
at infinity.
\end{slogan}
Let $X$ be a scheme. Let $f \in \Gamma(X, \mathcal{O}_X)$.
Denote $X_f \subset X$ the open where $f$ is invertible, see
Schemes, Lemma \href{https://stacks.math.columbia.edu/tag/01HZ}{lemma f open (Tag 01HZ)}.
If $X$ is quasi-compact and quasi-separated, the canonical map
$$
\Gamma(X, \mathcal{O}_X)_f \longrightarrow \Gamma(X_f, \mathcal{O}_X)
$$
is an isomorphism. Moreover, if $\mathcal{F}$ is a quasi-coherent
sheaf of $\mathcal{O}_X$-modules the map
$$
\Gamma(X, \mathcal{F})_f \longrightarrow \Gamma(X_f, \mathcal{F})
$$
is an isomorphism.
\end{lemma}

\begin{proof}
Write $R = \Gamma(X, \mathcal{O}_X)$. Consider the canonical morphism
$$
\varphi : X \longrightarrow \Spec(R)
$$
of schemes, see
Schemes, Lemma
\href{https://stacks.math.columbia.edu/tag/01I1}{lemma morphism into affine (Tag 01I1)}.
Then the inverse image of the standard open $D(f)$ on the
right hand side is $X_f$ on the left hand side.
Moreover, since $X$ is assumed quasi-compact and quasi-separated
the morphism $\varphi$ is quasi-compact and quasi-separated,
see Schemes, Lemma \href{https://stacks.math.columbia.edu/tag/01K4}{lemma quasi compact affine (Tag 01K4)} and
\href{https://stacks.math.columbia.edu/tag/01KV}{lemma compose after separated (Tag 01KV)}. Hence by
Schemes, Lemma \href{https://stacks.math.columbia.edu/tag/01LC}{lemma push forward quasi coherent (Tag 01LC)}
we see that $\varphi_*\mathcal{F}$ is quasi-coherent.
Hence we see that $\varphi_*\mathcal{F} = \widetilde M$
with $M = \Gamma(X, \mathcal{F})$ as an $R$-module.
Thus we see that
$$
\Gamma(X_f, \mathcal{F}) =
\Gamma(D(f), \varphi_*\mathcal{F}) =
\Gamma(D(f), \widetilde M) = M_f
$$
which is exactly the content of the lemma. The first displayed isomorphism
of the lemma follows by taking $\mathcal{F} = \mathcal{O}_X$.
\end{proof}


\begin{lemma}
\label{more-algebra-lemma-regular-koszul-regular}
Let $R$ be a ring, $M$ an $R$-module, and $f_1, \ldots, f_r \in R$
such that for $i = 1, \ldots, r$ multiplication by $f_i$
is injective on $M/(f_1, \ldots, f_{i - 1})M$. Then $f_1, \ldots, f_r$
is $M$-Koszul regular. In particular, an $M$-regular sequence
is $M$-Koszul-regular and any regular sequence is Koszul-regular.
\end{lemma}

\begin{proof}
Let $R$, $M$, $f_1, \ldots, f_r$ be as in the first sentence of the lemma.
If $r = 1$, it is immediate that $f_1$ is $M$-Koszul-regular.
Assume $r > 1$. Since $f_1$ is a nonzerodivisor on $M$, we obtain
a short exact sequence of complexes:
$$
0 \to K_\bullet(f_2, \ldots, f_r) \otimes M
\xrightarrow{f_1}
K_\bullet(f_2, \ldots, f_r) \otimes M \to
K_\bullet(\overline{f}_2, \ldots, \overline{f}_r) \otimes M/f_1M \to 0
$$
Here $\overline{f}_i$ is the image of $f_i$ in $R/(f_1)$. By
Lemma \href{https://stacks.math.columbia.edu/tag/0629}{lemma cone koszul (Tag 0629)}
the complex $K_\bullet(f_1, \ldots, f_r)$
is isomorphic to the cone of multiplication by $f_1$
on $K_\bullet(f_2, \ldots, f_r)$. Thus
$K_\bullet(R, f_1, \ldots, f_r) \otimes M$ is isomorphic
to the cone on the first map. Hence
$K_\bullet(\overline{f}_2, \ldots, \overline{f}_r) \otimes M/f_1M$
is quasi-isomorphic to $K_\bullet(f_1, \ldots, f_r) \otimes M$.
As $R/(f_1)$, $M/f_1M$, $\overline{f}_2, \ldots, \overline{f}_r$
satisfy the conditions of the lemma, by induction we conclude
this complex is acyclic in postive degrees. This finishes the proof
of the first statement. The second statement immediately follows
from the first.
\end{proof}


\begin{lemma}
\label{constructions-lemma-proj-sheaves}
Let $S$ be a graded ring. Let $M$ be a graded $S$-module.
Let $\widetilde M$ be the sheaf of $\mathcal{O}_{\text{Proj}(S)}$-modules
associated to $M$.
\begin{enumerate}
\item For every $f \in S$ homogeneous of positive degree we have
$$
\Gamma(D_{+}(f), \mathcal{O}_{\text{Proj}(S)}) = S_{(f)}.
$$
\item For every $f\in S$ homogeneous of positive degree
we have $\Gamma(D_{+}(f), \widetilde M) = M_{(f)}$
as an $S_{(f)}$-module.
\item Whenever $D_{+}(g) \subset D_{+}(f)$ the restriction mappings
on $\mathcal{O}_{\text{Proj}(S)}$ and $\widetilde M$
are the maps
$S_{(f)} \to S_{(g)}$ and $M_{(f)} \to M_{(g)}$ from Lemma
\href{https://stacks.math.columbia.edu/tag/01M4}{lemma standard open (Tag 01M4)}.
\item Let $\mathfrak p$ be a homogeneous prime of $S$ not containing
$S_{+}$, and let $x \in \text{Proj}(S)$
be the corresponding point. We have
$\mathcal{O}_{\text{Proj}(S), x} = S_{(\mathfrak p)}$.
\item Let $\mathfrak p$ be a homogeneous prime of $S$ not containing
$S_{+}$, and let $x \in \text{Proj}(S)$
be the corresponding point. We have $(\widetilde M)_x = M_{(\mathfrak p)}$
as an $S_{(\mathfrak p)}$-module.
\item
\label{constructions-item-map}
There is a canonical ring map
$
S_0 \longrightarrow \Gamma(\text{Proj}(S), \widetilde S)
$
and a canonical $S_0$-module map
$
M_0 \longrightarrow \Gamma(\text{Proj}(S), \widetilde M)
$
compatible with the descriptions of sections over standard opens
and stalks above.
\end{enumerate}
Moreover, all these identifications are functorial in the graded
$S$-module $M$. In particular, the functor $M \mapsto \widetilde M$
is an exact functor from the category of graded $S$-modules
to the category of $\mathcal{O}_{\text{Proj}(S)}$-modules.
\end{lemma}

\begin{proof}
Assertions (1) - (5) are clear from the discussion above.
We see (6) since there are canonical maps $M_0 \to M_{(f)}$,
$x \mapsto x/1$ compatible with the restriction maps
described in (3). The exactness of the functor $M \mapsto \widetilde M$
follows from the fact that the functor $M \mapsto M_{(\mathfrak p)}$
is exact (see Algebra, Lemma \href{https://stacks.math.columbia.edu/tag/00JR}{lemma proj prime (Tag 00JR)})
and the fact that exactness of short exact sequences
may be checked on stalks, see
Modules, Lemma \href{https://stacks.math.columbia.edu/tag/01AG}{lemma abelian (Tag 01AG)}.
\end{proof}


\begin{proposition}
\label{proposition-quasi-compact-affine-diagonal}
Let $X$ be a quasi-compact scheme with affine diagonal.
Then the functor (\ref{equation-compare})
$$
D(\QCoh(\mathcal{O}_X))
\longrightarrow
D_\QCoh(\mathcal{O}_X)
$$
is an equivalence with quasi-inverse given by $RQ_X$.
\end{proposition}

\begin{proof}
Let $U \subset X$ be an affine open. Then the morphism
$U \to X$ is affine by
Morphisms, Lemma \href{https://stacks.math.columbia.edu/tag/01SG}{lemma affine permanence (Tag 01SG)}.
Thus the assumption of Lemma \ref{lemma-argument-proves}
holds by Lemma \ref{lemma-affine-pushforward} and we win.
\end{proof}


\noindent
Recall that given a scheme $X$, an invertible sheaf $\mathcal{L}$
on $X$, and a sheaf of $\mathcal{O}_X$-modules $\mathcal{F}$
we get a graded ring
$\Gamma_*(X, \mathcal{L}) =
\bigoplus\nolimits_{n \geq 0} \Gamma(X, \mathcal{L}^{\otimes n})$
and a graded $\Gamma_*(X, \mathcal{L})$-module
$\Gamma_*(X, \mathcal{L}, \mathcal{F}) =
\bigoplus\nolimits_{n \in \mathbf{Z}}
\Gamma(X, \mathcal{F} \otimes_{\mathcal{O}_X} \mathcal{L}^{\otimes n})$
see Modules, Definition \ref{modules-definition-gamma-star}.
If we have moreover a section $s \in \Gamma(X, \mathcal{L})$, then
we obtain a map
\begin{equation}
\label{properties-equation-module-invert-s}
\Gamma_*(X, \mathcal{L}, \mathcal{F})_{(s)}
\longrightarrow
\Gamma(X_s, \mathcal{F}|_{X_s})
\end{equation}
which sends $t/s^n$ where
$t \in \Gamma(X, \mathcal{F} \otimes_{\mathcal{O}_X} \mathcal{L}^{\otimes n})$
to $t|_{X_s} \otimes s|_{X_s}^{-n}$. This makes sense
because $X_s \subset X$ is by definition the open over which


\begin{lemma}
\label{lemma-generator-P1}
Let $A$ be a ring. Let $X = \mathbf{P}^n_A$. Then
$$
E =
\mathcal{O}_X \oplus \mathcal{O}_X(-1) \oplus \ldots \oplus \mathcal{O}_X(-n)
$$
is a generator
(Derived Categories, Definition \ref{derived-definition-generators})
of $D_\QCoh(X)$.
\end{lemma}

\begin{proof}
Let $K \in D_\QCoh(\mathcal{O}_X)$. Assume
$\Hom(E, K[p]) = 0$ for all $p \in \mathbf{Z}$.
We have to show that $K = 0$.
By Derived Categories, Lemma
\ref{derived-lemma-right-orthogonal}
we see that $\Hom(E', K[p])$ is zero for all $E' \in \langle E \rangle$
and $p \in \mathbf{Z}$.
By Lemma \ref{lemma-construct-the-next-one}
applied with $a = -n - 1$
we see that $\mathcal{O}_X(-n - 1) \in \langle E \rangle$
because it is quasi-isomorphic to a finite complex
whose terms are finite direct sums of summands of $E$.
Repeating the argument with $a = -n - 2$ we see that
$\mathcal{O}_X(-n - 2) \in \langle E \rangle$.
Arguing by induction we find that $\mathcal{O}_X(-m) \in \langle E \rangle$
for all $m \geq 0$.
Since
$$
\Hom(\mathcal{O}_X(-m), K[p]) =
H^p(X, K \otimes_{\mathcal{O}_X}^\mathbf{L} \mathcal{O}_X(m)) =
H^p(X, K \otimes_{\mathcal{O}_X}^\mathbf{L} \mathcal{O}_X(1)^{\otimes m})
$$
we conclude that $K = 0$ by Lemma \ref{lemma-nonzero-some-cohomology}.
(This also uses that $\mathcal{O}_X(1)$ is an ample
invertible sheaf on $X$ which follows from
Properties, Lemma \ref{properties-lemma-open-in-proj-ample}.)
\end{proof}
\end{document}
